\section{Task-Specific Solver Catalogue}
\label{sec:task-solver-catalogue}

Choose the physical task in this section, then use its named executable.
Each entry gives the equations, update order, inputs, outputs, and accuracy
checks needed to interpret the calculation. The corresponding catalogue is
\path{notes/SOLVER_MODULES.md}. Separate entry points validate each task's
inputs and call shared numerical kernels.  The old
\code{venus_modular_solver*}, weighted-PIC runner names, and
\code{run_wpic_subhalo_gpu*} remain compatibility interfaces for older
inputs.  New work should use a task-specific name so that the executable
itself rejects an input that belongs to another equation or collision family.
A maintained implementation is built and checked. Each entry states its
demonstrated scientific scope and the further convergence studies needed for
a new parameter range.

\subsection{Scope and entry-point matrix}

\begin{longtable}{@{}>{\raggedright\arraybackslash}p{0.31\linewidth}
                        >{\raggedright\arraybackslash}p{0.39\linewidth}
                        >{\raggedright\arraybackslash}p{0.22\linewidth}@{}}
\toprule
Task & Executable & Numerical scope \\
\midrule
\endfirsthead
\toprule
Task & Executable & Numerical scope \\
\midrule
\endhead
Expanding VP WPIC & \path{venus_cpu_vlasov_poisson_wpic_expanding} & maintained collisionless production task \\
Subhalo WPIC & \path{venus_cpu_subhalo_wpic} & maintained calibrated task; ensemble tests required \\
Constant material SIDM WPIC & \path{venus_cpu_sidm_wpic_constant} & maintained task; late-time locality caveat \\
Velocity-dependent material SIDM WPIC & \path{venus_cpu_sidm_wpic_velocity_dependent} & research extension; formula tests only \\
Constant SphDSMC SIDM & \path{venus_cpu_sidm_sphdsmc_constant}, \path{_openmp} & maintained constant/isotropic task; optional paper-exact mode \\
Velocity-dependent SphDSMC SIDM & \path{venus_cpu_sidm_sphdsmc_velocity_dependent}, \path{_openmp} & research extension; formula tests only \\
Low-compactness EV WPIC & \path{venus_cpu_einstein_vlasov_wpic_low_compactness} & maintained comparison example \\
High-compactness EV WPIC & \path{venus_cpu_einstein_vlasov_wpic_black_hole} & research task; separate Schwarzschild and ingoing-EF kernels \\
Grid Einstein--Vlasov & \path{venus_cpu_einstein_vlasov_eulerian}, \path{_openmp}, \path{_mpi} & grid task; convergence required \\
Grid Vlasov--Poisson & \path{venus_cpu_vlasov_poisson_eulerian}, \path{_openmp}, \path{_mpi} & grid task; convergence required \\
Grid BP exact finite angle & \path{venus_cpu_boltzmann_poisson_exact}, \path{_openmp} & exact quadrature \\
Grid BP reversible finite angle & \path{venus_cpu_boltzmann_poisson_reversible}, \path{_openmp} & reversible quadrature \\
Grid BP sampled reversible & \path{venus_cpu_boltzmann_poisson_reversible_sampled}, \path{_openmp} & sampled reversible \\
Grid BP streamed reversible & \path{venus_cpu_boltzmann_poisson_reversible_streamed}, \path{_openmp} & streamed reversible \\
Grid BP finite-angle selector & \path{venus_cpu_boltzmann_poisson_finite_angle}, \path{_openmp} & compatibility selector \\
Expanding prompt-cusp BP & \path{venus_cpu_boltzmann_poisson_expanding_prompt_cusp}, \path{_openmp} & developing finite-angle application; finite-box caveat \\
Grid BP Fokker--Planck & \path{venus_cpu_boltzmann_poisson_fokker_planck}, \path{_openmp} & small-angle task \\
Grid BP Landau & \path{venus_cpu_boltzmann_poisson_landau}, \path{_openmp} & Landau-named task \\
Canonical orbit-phase BP & \path{venus_cpu_boltzmann_poisson_canonical_orbit}, \path{_openmp} & developing orbit-atlas task; convergence required \\
GPU VP WPIC & \path{venus_gpu_vlasov_poisson} and \path{_mpi} & restart-based task \\
GPU subhalo BP & \path{venus_gpu_subhalo_boltzmann_poisson} and \path{_mpi} & restart-based task \\
GPU material SIDM dispatcher & \path{venus_gpu_sidm_boltzmann_poisson} and \path{_mpi} & static-halo restart task \\
GPU constant SIDM & \path{venus_gpu_sidm_constant} and \path{_mpi} & locality/convergence required \\
GPU power-law SIDM & \path{venus_gpu_sidm_power_law} and \path{_mpi} & locality/convergence required \\
GPU Yukawa--Born SIDM & \path{venus_gpu_sidm_yukawa_born} and \path{_mpi} & angular/convergence checks required \\
\bottomrule
\end{longtable}

In this compact matrix, \path{_openmp} and \path{_mpi} are literal executable
suffixes, not separate programs.  For example, the complete SphDSMC names are
\path{venus_cpu_sidm_sphdsmc_constant} and
\path{venus_cpu_sidm_sphdsmc_constant_openmp}.  The CPU Makefile creates an
\path{_openmp} form of every grid/SphDSMC task but creates \path{_mpi} forms
only for collisionless Eulerian Einstein--Vlasov and Vlasov--Poisson.  The GPU
CMake build creates the six full names ending in \path{_mpi} only when its MPI
option is enabled.  Weighted-marker CPU task runners have no suffix variants.

An implemented task is limited to its documented equation, representation,
boundary/restart contract, and parameter domain.  Every quantitative run still
requires the convergence and locality studies named in its module entry.

\subsection{Task-main source layout}

Every preferred executable has a distinct task main; numerical engines remain
shared.  The CPU task sources are:
\begin{itemize}[itemsep=0.2em]
\item \path{cpu/src/apps/einstein_vlasov/eulerian.cpp},
  \path{cpu/src/apps/einstein_vlasov/wpic_low_compactness.cpp}, and
  \path{cpu/src/apps/einstein_vlasov/wpic_black_hole.cpp};
\item \path{cpu/src/apps/vlasov_poisson/eulerian.cpp} and
  \path{cpu/src/apps/vlasov_poisson/expanding_wpic.cpp};
\item \path{cpu/src/apps/subhalo/expanding_wpic.cpp};
\item \path{cpu/src/apps/sidm/sphdsmc_constant.cpp},
  \path{cpu/src/apps/sidm/sphdsmc_velocity_dependent.cpp},
  \path{cpu/src/apps/sidm/wpic_constant.cpp}, and
  \path{cpu/src/apps/sidm/wpic_velocity_dependent.cpp}; and
\item the nine files under \path{cpu/src/apps/boltzmann_poisson}:
  \path{fokker_planck.cpp}, \path{landau.cpp}, \path{finite_angle.cpp},
  \path{exact.cpp}, \path{reversible.cpp},
  \path{reversible_sampled.cpp}, \path{reversible_streamed.cpp}, and
  \path{expanding_prompt_cusp.cpp}, plus \path{canonical_orbit.cpp}.
\end{itemize}
Shared CPU WPIC engines live in \path{cpu/src/drivers}; grid, collision,
parallel, and I/O kernels live in \path{cpu/include/ev}.

The six GPU task mains are
\path{gpu/tools/venus_gpu_vlasov_poisson.cu},
\path{gpu/tools/venus_gpu_subhalo_boltzmann_poisson.cu},
\path{gpu/tools/venus_gpu_sidm_boltzmann_poisson.cu},
\path{gpu/tools/venus_gpu_sidm_constant.cu},
\path{gpu/tools/venus_gpu_sidm_power_law.cu}, and
\path{gpu/tools/venus_gpu_sidm_yukawa_born.cu}.  They call the common engine in
\path{gpu/src/wpic_production_driver.cu}; device algorithms remain under
\path{gpu/src} and public contracts under \path{gpu/include/venus_gpu}.
Thus the command and parser are task-specific while the numerical algorithm is
implemented once.

\subsection{CPU modules, one by one}
\subsubsection{Expanding Vlasov--Poisson WPIC}
\paragraph{Physics scope and exclusions.}
This marker task is explicitly collisionless; it does not
construct a material or subhalo collision clock.
\paragraph{Equations/algorithm.}
Weighted markers represent the matter distribution. The solver first
imports or generates them, deposits their spherical moments, and solves
Poisson's equation. It advances the markers by expanding KDK, applies the
boundaries, and writes profiles, histories, and restarts. Repeating this cycle
evolves collisionless VP without a collision scheduler.
\paragraph{Source/entry/lock.}
Use \code{venus_cpu_vlasov_poisson_wpic_expanding}.  Its task main is
\path{cpu/src/apps/vlasov_poisson/expanding_wpic.cpp}; the complete run loop,
IC generation, restart reader, and output writer are in
\path{cpu/src/drivers/expanding_wpic.cpp}, with marker algorithms in
\path{cpu/include/ev/weighted_vlasov_pic.hpp}.  Collision strength is locked
to zero and all material/subhalo controls are rejected.
\paragraph{Contract and limitations.}
Slope, marker/radial resolution, background, timestep, deposition, boundary,
cadence, IC/restart, and threads apply.  It writes profile snapshots, history,
metadata, mesh contract, and WPIC restart, in one shared-memory process.  Use
\path{cpu/tools/render_weighted_pic_pure_gamma_gifs.py} or the equilibrium renderer.
The supplied controls illustrate the run path; quantitative use requires
marker/cloud/background/boundary convergence.

\subsubsection{Subhalo collisional WPIC}
\paragraph{Physics scope and exclusions.}
This task treats each marker as a complete Plummer object
and models calibrated scatter or capture.  It is not material SIDM.
\paragraph{Equations/algorithm.}
Velocity/mass-dependent calibrated cross sections select mutually exclusive
scatter/capture.  WPIC first advances to the earliest field or collision deadline. Radial
cells define the local partners, and exact pair majorants bound their trial
rates. Counter-based random draws select and accept pair encounters. The physical
encounter law classifies scattering or capture. The solver
updates or merges the markers, records internal-energy compaction, and
redeposits the matter. It checks invariants, locality, and event intervals
before saving the clocks and random counters.
\paragraph{Source/entry/lock.}
Use \code{venus_cpu_subhalo_wpic}.  Its task main is
\path{cpu/src/apps/subhalo/expanding_wpic.cpp}; the run loop and I/O are in
\path{cpu/src/drivers/expanding_wpic.cpp}, and the collision algorithm is in
\path{cpu/include/ev/weighted_vlasov_pic_subhalo_collisions.hpp} plus the event scheduler.
Positive subhalo strength is required; material/grid/SphDSMC controls fail.
\paragraph{Contract and limitations.}
Subhalo table/strength, Plummer/low-speed, collision interval/fraction/CFL,
lanes/locality/seed, optional per-cell adaptive clock, marker/grid,
output/restart, and threads apply.  Outputs add a collision-history table.
Safety stops return documented nonzero process exit values; this CPU driver
does not write a separate terminal-status file.  Execution is shared-memory.
Use the subhalo comparison/time/age/strength GIFs and example case 04, whose
configured endpoint is \(t=50\).  Population conclusions require ensemble,
resolution, collision-cell, and capture-compaction studies.

\subsubsection{Constant material SIDM WPIC}
\paragraph{Physics scope and exclusions.}
This task samples microscopic constant \(\sigma/m\) SIDM with weighted
markers.  The packaged deep-collapse run has a stated late nonlocal-shell
warning; post-warning structure is not locality-resolved science.
\paragraph{Equations/algorithm.}
Pair area is \((m_i+m_j)(\sigma/m)/2\).  After importing markers, the solver deposits their matter and advances
KDK gravity to the next field or collision event. Adaptive shells define
local NTC partners, which scatter isotropically in their center-of-mass
frame. The field is reconstructed when required. Invariant, effective-number,
and MFP checks precede the checkpoint.
\paragraph{Source/entry/lock.}
Use \code{venus_cpu_sidm_wpic_constant}.  Its task main is
\path{cpu/src/apps/sidm/wpic_constant.cpp};
\path{cpu/src/drivers/sidm_wpic.cpp} owns the command line, complete loop, and
I/O, while \path{cpu/include/ev/weighted_vlasov_pic_collisions.hpp} owns the
collision steps.  Event/remap and Frank-import helpers are shared.
Constant material SIDM is required; velocity-dependent, subhalo, SphDSMC, and
grid controls are rejected.
\paragraph{Contract and limitations.}
\(\sigma_0\), strength, shell/CFL/substep/locality, marker/grid/cloud,
event/cadence/restart, thread controls, and the optional local shell clock and
outer interval cap apply.  The local clock requires constant cross section,
equal shell masses, and CFL no larger than one.  Exactly one of
\code{--frank-restart} and \code{--pic-restart} is required.  The driver writes
\path{profile_NNNNNNNN.csv}, where the eight digits encode the step,
\path{history.csv},
\path{operator_mass_budget.csv}, and atomic \path{restart_state.bin}; it does
not write \path{metadata.json}.  It uses one shared-memory process.  Use the
Frank WPIC renderers and example case 05.
Post-warning science requires a finer locality-converged rerun.

\subsubsection{Velocity-dependent material SIDM WPIC}
\paragraph{Physics scope and exclusions.}
This research task locks material WPIC to power-law or Yukawa--Born velocity
dependence; no supplied long-time example establishes core evolution.
\paragraph{Equations/algorithm.}
The module above is retained, but \(\Sigma_{ij}(u)\) changes the pair majorant
and Yukawa--Born draws a forward-peaked differential angle.  This CPU path
retains the mesh-wide reference collision clock.
\paragraph{Source/entry/lock.}
Use \code{venus_cpu_sidm_wpic_velocity_dependent}.  Its task main is
\path{cpu/src/apps/sidm/wpic_velocity_dependent.cpp}; it uses the same
\path{cpu/src/drivers/sidm_wpic.cpp} driver and
\path{cpu/include/ev/weighted_vlasov_pic_collisions.hpp} algorithms as the
constant task.  \code{powerlaw} or \code{yukawa_born} is mandatory.  Required
scale/exponent or coupling/masses are checked; isotropic Yukawa, constant,
subhalo, SphDSMC, and grid inputs fail.
\paragraph{Contract and limitations.}
IC/output/shared-memory/plotting match constant WPIC.  Formula, angular,
majorant, and invariant agreement does not establish long-time relaxation.
Compare each law with a constant cross section giving the same initial
relaxation time. Check Maxwellian relaxation and core evolution, then refine
collision-shell locality and marker number across independent realizations.

\subsubsection{Constant material SIDM with SphDSMC}
\paragraph{Physics scope and exclusions.}
This maintained path implements persistent-marker, constant/isotropic
material SIDM.  The ordinary task profile does not by itself claim the
published paper-exact parameter set.  That stricter mode is enabled only by
\code{EV_SPHERICAL_DSMC_PAPER_EXACT=1}; it checks the published shell size,
shell floor, collision step, gravity step, orbit retry, and energy tolerance.
The supplied example uses 512 target particles per collision shell rather
than the paper-exact value 4096, and it is not a completed gravothermal
catastrophe.
\paragraph{Equations/algorithm.}
Local exposure is proportional to \(u\Sigma_{ij}\Delta t/V_{\rm shell}\),
with constant \(\sigma/m\), and elastic COM scattering closes pair invariants.
The solver sorts markers by radius so enclosed mass determines the force.
That force sets the orbit timestep, and stiff trajectories can be retried
with finer integration. Adaptive shells then define local NTC partners and
isotropic scattering angles. Invariant and MFP-locality checks precede output
and restart writing.
\paragraph{Source/entry/lock.}
Use \code{venus_cpu_sidm_sphdsmc_constant} for the portable threaded build or
\code{venus_cpu_sidm_sphdsmc_constant_openmp} for the OpenMP build.  Their
task main is \path{cpu/src/apps/sidm/sphdsmc_constant.cpp}; implementation is in
\path{cpu/include/ev/spherical_dsmc_solver.hpp} and
\path{cpu/include/ev/spherical_dsmc_collision.hpp}.  BP+SphDSMC+constant/isotropic is
required; every velocity-dependent, grid, WPIC, and subhalo mode is rejected.
\paragraph{Contract and limitations.}
Particle mass/count, \(\sigma_0\), shell occupancy, CFL, locality,
paper-exact, orbit, boundary, cadence, IC/restart, and threads apply.  It
samples a distribution function or resumes a persistent-marker restart and writes metadata,
profiles, history, mass budget, and restart.  Run exactly one operating-system
process: the unsuffixed executable uses portable C++ threads and the
\path{_openmp} executable uses OpenMP threads.  There is no SphDSMC MPI
executable.
Use \path{cpu/tools/render_frank_venus_sphdsmc_gif.py} and example case 06.  The
finite-time example does not complete gravothermal collapse; a quantitative
calculation requires full-trajectory resolution, locality, and
thread-equivalence studies.

\subsubsection{Velocity-dependent material SIDM with SphDSMC}
\paragraph{Physics scope and exclusions.}
This research path exposes power-law, softened-gravity, and Yukawa--Born
SphDSMC without calling them paper-exact constant SIDM.
\paragraph{Equations/algorithm.}
The exposure uses \(u\Sigma_{ij}(u)\); Yukawa--Born also draws the angle from
the differential law.  It follows the seven steps above with a
velocity-dependent majorant and law-specific angular sampler.
\paragraph{Source/entry/lock.}
Use \code{venus_cpu_sidm_sphdsmc_velocity_dependent} or the complete OpenMP
name \code{venus_cpu_sidm_sphdsmc_velocity_dependent_openmp}.  Their task main
is \path{cpu/src/apps/sidm/sphdsmc_velocity_dependent.cpp}; the shared
implementation is \path{cpu/include/ev/spherical_dsmc_solver.hpp} and
\path{cpu/include/ev/spherical_dsmc_collision.hpp}.  The task requires a
nonconstant law, rejects paper-exact constant mode, and validates all required
scale/exponent or Yukawa parameters.  An isotropic Yukawa alias is invalid.
\paragraph{Contract and limitations.}
IC/restart, SphDSMC output, one-process portable-thread/OpenMP execution, and
the SphDSMC renderer match the constant task.  No supplied long-time example
establishes core evolution.
Each law requires separate Maxwellian, rate, mean-free-path, resolution, and
core-evolution studies.

\subsubsection{Low-compactness Einstein--Vlasov WPIC}
\paragraph{Physics scope and exclusions.}
This weighted-marker task compares relativistic EV and Newtonian VP from the
same weighted realization; it is not the Eulerian solver.
\paragraph{Equations/algorithm.}
Invariant marker weights supply the relativistic matter moments. Each step
uses these moments to solve the constraints and advance the EV markers,
alongside Newtonian VP markers with the same sampled positions and momenta.
Each module assigns the weights appropriate to its phase-space measure.
Paired diagnostics and rotating restarts preserve the comparison.
\paragraph{Source/entry/lock.}
Use \code{venus_cpu_einstein_vlasov_wpic_low_compactness}.  Its task main is
\path{cpu/src/apps/einstein_vlasov/wpic_low_compactness.cpp}; the complete run
loop and paired restart are in
\path{cpu/src/drivers/einstein_vlasov_wpic.cpp}, with shared physics in
\path{cpu/include/ev/weighted_einstein_vlasov_pic.hpp}.  The low scenario is fixed and
high-shell/horizon/collision controls are rejected.
\paragraph{Contract and limitations.}
Low marker/radial resolution, mass, velocity, importance, CFL, endpoint,
cadence, and threads apply.  Output/restart carries paired EV/VP fields.
Only shared-memory execution is supported.  Use the low branch of
\path{cpu/tools/render_weighted_einstein_vlasov_pic_profile_gifs.py} and example case
01.  The supplied resolution is a virialized-halo example; a precision EV/VP
density comparison requires marker/radial/time refinement.

\subsubsection{High-compactness Einstein--Vlasov WPIC}
\paragraph{Physics scope and exclusions.}
This neutral spherical task follows a Plummer sphere or infalling shell in
either a Schwarzschild chart or an ingoing
Eddington--Finkelstein/advanced-null chart.  Particle angular momentum is an
isotropic magnitude, not net spacetime spin; the neutral black-hole exterior
is Schwarzschild, not Kerr.  Charge and Maxwell fields are absent.
\paragraph{Equations/algorithm.}
The Schwarzschild kernel stores orthonormal \(w\), solves its original radial
constraints, and stops before \(1-2m/r=0\).  The independent advanced-null
kernel stores canonical \(P=p_r>0\), uses
\(\dd s^2=-ab^2\dd v^2+2b\dd v\dd r+r^2\dd\Omega^2\), deposits \(T_{rr}\)
and \(S\), solves for \(b\) and \(mb\), and locates an apparent horizon at
\(a=1-2m/r=0\).  It advances \(\log P\), preserving \(P>0\) at every accepted
RK stage, and can continue with causally buffered trapped-region excision.
\paragraph{Source/entry/lock.}
Use \code{venus_cpu_einstein_vlasov_wpic_black_hole} with exactly one of
\code{--coordinates schwarzschild} and
\code{--coordinates eddington_finkelstein}.  Its task main is
\path{cpu/src/apps/einstein_vlasov/wpic_black_hole.cpp}, and its complete
command line, run loop, restart, and output logic are in
\path{cpu/src/drivers/einstein_vlasov_wpic.cpp}.  The separate implementations are
\path{cpu/include/ev/weighted_einstein_vlasov_pic_schwarzschild.hpp} and
\path{cpu/include/ev/weighted_einstein_vlasov_pic_eddington_finkelstein.hpp}; the facade
\path{cpu/include/ev/weighted_einstein_vlasov_pic.hpp} exposes both.  The runner fixes the
high scenario and rejects low-comparison and collision options.
\paragraph{Contract and limitations.}
High IC/mass/geometry/infall, grid/boundary/gap, resolution/cadence/thread
controls apply, with fixed or scaled particle angular momentum and EF-only
excision controls.  The coordinate-neutral physical-IC file fixes the common
physical family.  A fresh EF run additionally requires its solved-null
initial-data artifact.  Dynamical \(t=\mathrm{const}\) and \(v=\mathrm{const}\) surfaces
are not the same hypersurface; compare invariant mass/compactness milestones,
not equal coordinate times.  The EF chart has no origin-reflection semantics
and fails closed if an active characteristic reaches \(r=0\).  Shared-memory
execution and coordinate-specific histories/profiles/restarts are supported.
Use the high renderer and the Einstein--Vlasov analyzer.  Quantitative claims require
marker/radial/timestep/domain/excision convergence, constraint residuals, a
static transformed-solution test, pre-horizon cross-coordinate checks, and a
resolution-repeated angular-momentum bracket.
Section~\ref{sec:ingoing-eddington-finkelstein-wpic} gives every field,
source, characteristic, coordinate map, discretization, and failure condition.

\subsubsection{Eulerian Einstein--Vlasov}
\paragraph{Physics scope and exclusions.}
This grid task evolves a resolved spherical
\((r,w,L)\) distribution and metric.  It is neither the weighted-marker
low/high task nor a black-hole finder.
\paragraph{Equations and algorithm.}
For \(\dd s^2=-e^{2\mu}\dd\tau^2+e^{2\lambda}\dd r^2+r^2\dd\Omega^2\), \(f\) is constant
on relativistic characteristics and its density, current, and pressures source
the spherical constraints.  The numbered cycle is: (1) read/construct \(f\);
(2) integrate moments; (3) solve \(\lambda,\mu,m\); (4) select the occupied-orbit
timestep; (5) trace arrival points backward; (6) interpolate with physical
boundaries; (7) recompute constraints; (8) output/restart.  Full equations are
in the Einstein--Vlasov chapter.
\paragraph{Source, executable, and lock.}
Use \code{venus_cpu_einstein_vlasov_eulerian}, \code{_openmp}, or
\code{_mpi}.  The task main is
\path{cpu/src/apps/einstein_vlasov/eulerian.cpp};
\path{cpu/include/ev/solver.hpp} owns the complete time-step loop, and shared
numerical sources are
\path{cpu/include/ev/moments.hpp}, \path{cpu/include/ev/metric.hpp},
\path{cpu/include/ev/characteristic_integrators.hpp},
\path{cpu/include/ev/advection.hpp}, \path{cpu/include/ev/timestep.hpp}, and the common solver loop.  The
front end requires \code{physics_model=einstein_vlasov} and rejects VP/BP and
collision settings.
\paragraph{Run contract.}
Grid, metric/compactness, characteristic, boundary, IC, timestep, output,
restart, thread, and supported MPI controls apply; collision parameters do
not.  Analytic or \code{binary_rwl} ICs lead to metadata, history, profiles,
phase snapshots, and versioned restarts.  Serial, OpenMP, and supported MPI
\(L\)-slabs are available.  Plot with
\path{cpu/tools/plot_einstein_vlasov_outputs.py}; use the two Eulerian EV example INIs.
The supplied INIs are controlled starting problems, not long-time convergence
results.
Quantitative use requires constraint, resolution, boundary, restart, and
parallel-equivalence studies plus a trapped-surface diagnostic for horizon
claims.

\subsubsection{Eulerian Vlasov--Poisson}
\paragraph{Physics scope and exclusions.}
This grid task covers deterministic static/expanding VP with nodal,
finite-volume, conservative SLDG, and Hamiltonian-aligned variants.  It is not
a marker-noise or collision calculation.
\paragraph{Equations and algorithm.}
In static variables,
\(\partial_t f+w\partial_r f+(L^2/r^3-\partial_r\Phi)\partial_w f=0\),
\(\partial_r\Phi=GM(<r)/r^2\), and
\(\partial_rM=4\pi r^2\rho\); the expanding form is derived later.  The cycle
is: (1) integrate density; (2) reconstruct mass/potential; (3) select the
transport map/flux; (4) advance with boundaries and optional midpoint field;
(5) update carrier/flux ledgers; (6) output/restart.
\paragraph{Source/entry/lock.}
Use \code{venus_cpu_vlasov_poisson_eulerian}, \code{_openmp}, or
\code{_mpi}.  The task main is
\path{cpu/src/apps/vlasov_poisson/eulerian.cpp};
\path{cpu/include/ev/solver.hpp} owns the complete loop.  Shared files include
\path{cpu/include/ev/vlasov_poisson.hpp},
\path{cpu/include/ev/conservative_sldg_transport.hpp},
\path{cpu/include/ev/hamiltonian_sldg_transport.hpp},
\path{cpu/include/ev/vp_expanding_interpolation.hpp}, and
\path{cpu/include/ev/delayed_inner_orbit_carrier.hpp}.  The task requires VP and no BP
operator; incompatible inflow/expanding/integrator combinations fail closed.
\paragraph{Contract and limitations.}
Universe, Poisson source, integrator, SLDG, boundary, regularization, grid, IC,
cadence, and parallel controls apply; metric/collision controls do not.  It
uses grid IC/restarts and grid outputs/ledgers.  Serial, OpenMP, and
collisionless MPI slabs are supported.  Use
\path{cpu/tools/spherical_profile_rendering.py} for named-column profile tables or the
Python \path{cpu/python/ev_solver/plotting.py} helpers for full grid output;
example run 03 and the VP configurations under
\path{cpu/examples/modular} provide inputs.  Finite-volume and SLDG results
require method-specific equilibrium, boundary, restart, MPI, and convergence
studies.

\subsubsection{Grid Boltzmann--Poisson: finite-angle family}
\paragraph{Physics scope and exclusions.}
This grid task exposes exact-quadrature, reversible, sampled, and
streamed finite-angle Eulerian operators.  It is not SphDSMC or WPIC.
\paragraph{Equations/algorithm.}
\(\partial_t f+\{f,H\}=C_{\rm B}[f]\), with
\(C_{\rm B}=\int\dd^3v_*\dd\Omega\,u(\dd\Sigma/\dd\Omega)
[f'f_*'-ff_*]\).  After transport, the solver assembles radial collision cells and their
quadrature nodes or sampled proposals. It evaluates loss and gain or
reversible increments, using collision substeps or streaming to satisfy
majorant and positivity limits. Outgoing states are remapped to the grid, and
mass, momentum, energy, omitted contributions, and sampling errors are
recorded in the corresponding ledgers.
\paragraph{Source/entry/lock.}
Use the backend-specific primary binary
\path{venus_cpu_boltzmann_poisson_exact},
\path{venus_cpu_boltzmann_poisson_reversible},
\path{venus_cpu_boltzmann_poisson_reversible_sampled}, or
\path{venus_cpu_boltzmann_poisson_reversible_streamed}, with \path{_openmp}
for the shared-memory build.  The broad
\path{venus_cpu_boltzmann_poisson_finite_angle} name is a compatibility
selector.  The task mains are the correspondingly named files
\path{cpu/src/apps/boltzmann_poisson/exact.cpp},
\path{cpu/src/apps/boltzmann_poisson/reversible.cpp},
\path{cpu/src/apps/boltzmann_poisson/reversible_sampled.cpp},
\path{cpu/src/apps/boltzmann_poisson/reversible_streamed.cpp}, and
\path{cpu/src/apps/boltzmann_poisson/finite_angle.cpp}.
\path{cpu/include/ev/solver.hpp} owns the full transport--collision loop;
operator files are
\path{cpu/include/ev/finite_angle_boltzmann.hpp},
\path{cpu/include/ev/reversible_finite_angle_boltzmann.hpp},
\path{cpu/include/ev/sampled_reversible_boltzmann.hpp}, and
\path{cpu/include/ev/streamed_reversible_boltzmann.hpp}.  Every primary binary locks exactly
one operator and rejects FP, SphDSMC, and every other finite-angle backend.
\paragraph{Contract and limitations.}
The selected differential/total law, quadrature/proposal, majorant,
replicates/workspace, cadence, and grid/VP controls apply; FP flux/RKL settings
do not select this operator.  Yukawa total and differential kernels must be
consistent.  It uses grid IC/restart/output plus operator ledgers in one
process.  The unsuffixed task uses portable C++ threads; the
\path{_openmp} task uses OpenMP threads.  No finite-angle MPI task exists.  Use
\path{cpu/tools/spherical_profile_rendering.py}.  Each submethod still
needs its own equilibrium,
quadrature/proposal, restart, and long-time convergence gate.

\subsubsection{Expanding prompt-cusp Boltzmann--Poisson}
\paragraph{Physics scope and exclusions.}
This task locks the expanding finite-angle setup intended to form a
prompt-cusp-like inner profile and an NFW-like outer transition.  One finite
box does not establish an asymptotic \(r^{-3}\) tail.
\paragraph{Equation/numbered algorithm.}
It composes the expanding VP equations and finite-angle \(C_{\rm B}\) above:
(1) read the \(\gamma_i\) IC; (2) expanding transport; (3) apply the registered
collision cadence/current minimum quadrature; (4) close ledgers; (5) write
profiles/restart through the endpoint.
\paragraph{Source/entry/lock.}
Use \code{venus_cpu_boltzmann_poisson_expanding_prompt_cusp}; it requires
expanding BP, finite-angle collisions, and nonzero strength, and rejects
static, collisionless, FP, and SphDSMC inputs.
The task main is
\path{cpu/src/apps/boltzmann_poisson/expanding_prompt_cusp.cpp}; the complete
loop is \path{cpu/include/ev/solver.hpp}, and the finite-angle operator is
\path{cpu/include/ev/finite_angle_boltzmann.hpp}.
\paragraph{Contract and limitations.}
The slope, expansion, box/grid, kernel/quadrature, strength, cadence, IC, and
output controls apply.  The operator requires \(4\times2\times4\) or higher
quadrature.  The portable successful-run \path{params.ini} already uses that
minimum.  It uses a phase-grid IC/restart and grid BP output in one process.
The unsuffixed task uses portable C++ threads; the \path{_openmp} task uses
OpenMP threads.  Plot with \path{cpu/tools/spherical_profile_rendering.py}.
Quantitative profile interpretation requires a grid/box/quadrature hierarchy.

\subsubsection{Grid Boltzmann--Poisson: Fokker--Planck/Landau}
\paragraph{Physics scope and exclusions.}
This grid task adds a small-angle collision approximation to Eulerian
VP; it is neither finite-angle DSMC nor subhalo capture.
\paragraph{Equations/algorithm.}
The collision term is
\(C_{\rm FP}=-\nabla_v\!\cdot(\bm A f)+\tfrac12\partial_{v_i}\partial_{v_j}
(D_{ij}f)\).  After VP transport, local collision moments determine the drift and
diffusion coefficients. Projecting those coefficients onto velocity faces
gives conservative fluxes. The solver advances these fluxes explicitly or by
RKL2 and then checks positivity and the conservation ledgers.
\paragraph{Source/entry/lock.}
Use \code{venus_cpu_boltzmann_poisson_fokker_planck} (or \code{_openmp}) for
the Fokker--Planck name and \code{venus_cpu_boltzmann_poisson_landau} (or
\code{_openmp}) for the Landau-named task.  Their task mains are
\path{cpu/src/apps/boltzmann_poisson/fokker_planck.cpp} and
\path{cpu/src/apps/boltzmann_poisson/landau.cpp};
\path{cpu/include/ev/solver.hpp} owns the complete loop, and
\path{cpu/include/ev/boltzmann_poisson.hpp} owns operator dispatch.  Each
binary locks its accepted operator token and rejects
finite-angle, SphDSMC, and canonical choices.
\paragraph{Contract and limitations.}
FP bath, coefficient refresh, collision flux/integrator, RKL2, cross section,
strength, cadence, grid, and VP controls apply; finite-angle quadrature,
pair-majorant, and marker controls do not.  It uses grid IC/restart and adds
coefficient/flux/invariant ledgers.  Run one process: the unsuffixed task uses
portable C++ threads and the \path{_openmp} task uses OpenMP threads.  There is
no Fokker--Planck/Landau MPI executable.  Use
\path{cpu/tools/render_density_slope_sigma2_gif.py} and
\path{cpu/tools/analyze_sidm_benchmark.py}.  Quantitative use requires
analytic-relaxation and finite-angle-reference convergence.  Distributed
collision exchange is not implemented.

\subsubsection{Canonical orbit-phase Boltzmann--Poisson}
\paragraph{Physics scope and exclusions.}
This CPU task evolves a distribution on a bound-orbit atlas in canonical phase
and couples transport to local finite-angle collision events and a
self-consistent Poisson field.  It is not the Eulerian \((r,w,L)\) solver,
SphDSMC, or the weighted-marker WPIC event library.
\paragraph{Equations/algorithm.}
Canonical BP translates a distribution on a bound-orbit phase atlas, lifts
finite-angle gain/loss into orbit cells/open packets, compresses positive
packets, and rebuilds Poisson at events.  One update transports the atlas,
selects local NTC pairs, maps outgoing states to bound cells or open packets,
closes invariant ledgers, compresses packets, reconstructs the field, and
advances event clocks.
\paragraph{Implementation, executable, and lock.}
Use \path{venus_cpu_boltzmann_poisson_canonical_orbit} or the complete OpenMP
name \path{venus_cpu_boltzmann_poisson_canonical_orbit_openmp}; their main is
\path{cpu/src/apps/boltzmann_poisson/canonical_orbit.cpp}.  Shared
implementation lives in the canonical headers under \path{cpu/include/ev}.
The executable uses
its orbit-atlas command-line schema and does not accept a modular INI or a
WPIC/SphDSMC task configuration.
\paragraph{Contract, I/O, and parallelism.}
\code{--restart PATH} is always required; the implementation reads it as a
versioned Eulerian grid restart.  A stationary Eddington table is also always
required, either through \code{--eddington-table PATH} or through the default
\path{frank_core_eddington_8192.csv} beside that source restart.
\code{--canonical-restart PATH} optionally restores the evolved atlas and
event state, but does not replace either source input.  Controls cover
atlas/phase/radial resolution, transport and collision times, NTC
marker/shell sampling, event clocks, CFL/locality, open-packet compression,
reversible sampled/streamed operators, cadence, and threads.
Output includes \path{history.csv}, \path{profiles.csv},
\path{operator_mass_budget.csv}, and atomic \path{restart_state.bin}.  The task
uses one shared-memory process and schema-compatible profile/invariant tools.
\paragraph{Required checks.}
Establish atlas, phase, radial, event-timestep, collision-marker, shell,
open-packet, boundary, restart, and thread-count convergence.  The integrated
\path{cpu/include/ev/weighted_pic_boltzmann_event.hpp} library remains an internal algorithm
and API rather than a separate public solver task.

\subsubsection{Python reference package}
\path{cpu/python/ev_solver} separates configuration, grids, ICs, moments,
metric, numerics, state, setup/evolution, I/O, and plotting for readable
reference work.  It is not a high-performance backend.  Compare its results
explicitly with a C++ task when using it as a numerical reference.

\subsection{GPU modules, one by one}

The CUDA implementation has its own state ownership, kernels, restart path,
and MPI layer.  GPU IC generation is host-only by design.

\subsubsection{GPU Vlasov--Poisson WPIC}
\paragraph{Physics scope and exclusions.}
This restart-based collisionless task is not EV, Eulerian
VP, canonical BP, or SphDSMC.
\paragraph{Equations/algorithm.}
It solves spherical VP with CIC/TSC deposition, enclosed-mass Poisson, and
static/expanding KDK: (1) read/upload restart; (2) deposit; (3) scan mass and
force; (4) fixed-step KDK, shortened only at output/final time; (5) boundaries;
(6) requested output/download; (7) repeat without collision scheduling.
\paragraph{Source/entry/lock.}
Use \code{venus_gpu_vlasov_poisson} or
\code{venus_gpu_vlasov_poisson_mpi}.  Both build the task main
\path{gpu/tools/venus_gpu_vlasov_poisson.cu}; the shared command line, run
loop, restart, and output logic are in
\path{gpu/src/wpic_production_driver.cu}.  State, VP, and MPI code live in
\path{gpu/src/wpic_state.cu}, \path{gpu/src/wpic_vlasov_poisson.cu}, and
\path{gpu/src/wpic_mpi.cpp}.  Collision strength is fixed to zero; collision controls
and unsupported delayed-marker restarts fail. Single-GPU VP accepts
static total-density or EdS density-contrast evolution, with EdS as its default.
MPI VP retains EdS. Both single-GPU backgrounds use the fixed KDK cap with
output/final-time shortening.
\paragraph{Contract and limitations.}
Restart/output/endpoint/timestep/mesh/cadence and supported background/
boundary state apply.  It writes \path{profile_NNNNN.csv},
\path{history_gpu.csv}, \path{metadata.json},
\path{radial_mesh_contract.txt}, and \path{restart_state.bin}; a successful
endpoint creates \path{completed.ok}.  One GPU or one rank/GPU is
supported.  Use common WPIC renderers or
\path{gpu/scripts/render_venus_gpu_density_slope_gif.py}.
\code{benchmark_vlasov_poisson_gpu} is the focused performance program.
Delayed-inner and broader boundary modes are absent.

\subsubsection{GPU subhalo Boltzmann--Poisson}
\paragraph{Physics scope and exclusions.}
This restart task combines GPU WPIC with calibrated
subhalo scatter/capture, not material SIDM.
\paragraph{Equations/algorithm.}
It uses the CPU subhalo WPIC task's calibrated pair areas and
COM/capture energy map:
(1) KDK to deadline; (2) cached exact pair majorants; (3) deterministic
trials; (4) scatter/capture; (5) compaction/migration; (6) redeposit;
(7) clock/locality/invariant audit; (8) gathered checkpoint.  Reference mode
uses one mesh-wide collision step; adaptive mode gives each cell an independent
CFL clock, batches local GPU steps, and stops a cell after capture.
\paragraph{Source/entry/lock.}
Use \code{venus_gpu_subhalo_boltzmann_poisson} or the complete distributed
name \code{venus_gpu_subhalo_boltzmann_poisson_mpi}.  Both build the task main
\path{gpu/tools/venus_gpu_subhalo_boltzmann_poisson.cu}; the shared run loop is
\path{gpu/src/wpic_production_driver.cu}, and collision code is
\path{gpu/src/wpic_subhalo_collisions.cu} together with VP, state, restart, and
MPI code.  Positive
subhalo strength is required and material/collisionless inputs fail.
\paragraph{Contract and limitations.}
Subhalo, mesh/cadence, local-clock/kernel-batch, restart/output, and device/MPI controls apply; material
\(\sigma/m\) controls do not.  Outputs add
\path{collision_history_gpu.csv}; a successful endpoint creates
\path{completed.ok}.  One GPU and one-rank/GPU ownership/migration are in the
supported scope.  Use the supplied GPU subhalo restart example and subhalo
renderers.  Quantitative use requires ensembles, rank scaling, and long-time
statistical convergence.

\subsubsection{GPU material SIDM Boltzmann--Poisson dispatcher}
\paragraph{Physics scope and exclusions.}
This static-halo task combines GPU WPIC and one material SIDM law.  The
material MPI scheduler is not rank-count invariant.
\paragraph{Equations/algorithm.}
Exposure is \(u\Sigma_{ij}(u)\) with elastic COM scattering: (1) upload
restart; (2) VP KDK to deadline; (3) fixed-occupancy local shells; (4) material
majorant; (5) pair/angle sampling; (6) device update; (7) invariant/collision
history; (8) clocks/RNG checkpoint.  Reference mode uses a mesh-wide collision
step; adaptive mode gives every fixed-occupancy shell its own CFL clock and
refreshes its velocity majorant after each local step.
\paragraph{Source/entry/lock.}
Use \code{venus_gpu_sidm_boltzmann_poisson} or the complete distributed name
\code{venus_gpu_sidm_boltzmann_poisson_mpi}.  Both build the task main
\path{gpu/tools/venus_gpu_sidm_boltzmann_poisson.cu}; the shared run loop is
\path{gpu/src/wpic_production_driver.cu}, and implementation is
\path{gpu/src/wpic_collisions.cu} plus shared VP/state/restart.  Static material SIDM
with one explicit law is required; subhalo capture, collisionless mode, and
unsupported expanding material inputs fail closed.
\paragraph{Contract and limitations.}
Restart/output/mesh/cadence, law, shell/CFL/interval/local-clock/RNG, and GPU/MPI controls
apply; subhalo table/compaction controls do not.  GPU WPIC output gains
\path{collision_history_gpu.csv}; a successful endpoint creates
\path{completed.ok}.  Single-GPU execution is supported;
material MPI results require rank-count and ownership studies.  Use the generic
SIDM analyzer/renderer.  \code{benchmark_material_sidm_gpu} is the focused
performance program.  Quantitative use requires CPU/GPU trajectories with the same constant
cross section and initial particles, followed by locality and convergence
studies.

\subsubsection{GPU constant material SIDM}
\paragraph{Scope/equation/algorithm.}
\code{venus_gpu_sidm_constant} and \code{_mpi} lock the preceding dispatcher
to \(\sigma(u)/m=\sigma_0\) and isotropic scattering; the generic eight-step
material algorithm is unchanged.
\paragraph{Lock and contract.}
The complete distributed name is \code{venus_gpu_sidm_constant_mpi}.  Both
executables build \path{gpu/tools/venus_gpu_sidm_constant.cu} and call
\path{gpu/src/wpic_production_driver.cu} and
\path{gpu/src/wpic_collisions.cu}.
\(\sigma_0\), strength, shell/event, restart/output, and GPU layout apply;
power-law/Yukawa/subhalo options fail.  GPU WPIC I/O and the generic SIDM plotting tools apply.
\paragraph{Required checks.}
Local formula, rate, and pair-invariant agreement does not establish a
quantitative trajectory.  Compare constant-SIDM trajectories from the same initial particles and
physical controls. Then test locality and convergence with marker number,
timestep, random seed, and MPI rank count.

\subsubsection{GPU power-law material SIDM}
\paragraph{Scope/equation/algorithm.}
\code{venus_gpu_sidm_power_law} and \code{_mpi} lock the material runner to
the softened power-law \(\sigma(u)/m\); the speed-dependent law enters the
majorant and scattering remains isotropic.
\paragraph{Lock and contract.}
The complete distributed name is \code{venus_gpu_sidm_power_law_mpi}.  Both
executables build \path{gpu/tools/venus_gpu_sidm_power_law.cu} and call
\path{gpu/src/wpic_production_driver.cu} and
\path{gpu/src/wpic_collisions.cu}.
Normalization, reference speed, exponent/softening, strength, shell/event,
restart/output, and GPU layout apply.  Constant-only, Yukawa-differential, and
subhalo controls fail.  Common material I/O/plotting applies.
\paragraph{Required checks.}
Compare ensembles with constant and power-law cross sections that give the
same initial relaxation time. Then test locality and convergence with marker
number, timestep, random seed, and MPI rank count.

\subsubsection{GPU Yukawa--Born material SIDM}
\paragraph{Scope/equation/algorithm.}
\code{venus_gpu_sidm_yukawa_born} and \code{_mpi} lock the material runner to
the Born total cross section and normalized forward-peaked differential
angular CDF.
\paragraph{Lock and contract.}
The complete distributed name is \code{venus_gpu_sidm_yukawa_born_mpi}.  Both
executables build \path{gpu/tools/venus_gpu_sidm_yukawa_born.cu} and call
\path{gpu/src/wpic_production_driver.cu} and
\path{gpu/src/wpic_collisions.cu}.
Coupling, dark-particle/mediator masses, strength, shell/event, restart/output,
and GPU layout apply.  Isotropic angular aliases, power-law, and subhalo
controls fail.  Common material I/O/plotting applies.
\paragraph{Required checks.}
Establish forward-scattering resolution, rate, statistical CPU/GPU, locality,
and core-evolution convergence studies.

\subsubsection{GPU compatibility, restart, MPI, and benchmarks}
\code{run_wpic_subhalo_gpu} and \code{_mpi} remain auto-mode compatibility
names and must report their resolved task.  \path{gpu/src/wpic_restart.cpp} owns the
portable checkpoint and \path{gpu/src/wpic_mpi.cpp} owns one-rank/GPU
synchronization.  \code{benchmark_vlasov_poisson_gpu} and
\code{benchmark_material_sidm_gpu} are focused performance tasks;
\code{benchmark_wpic} is the compatibility aggregate.  Timing is not physics
convergence.  GPU EV, Eulerian VP/BP, canonical BP, SphDSMC, delayed-inner
return, and material Plummer turning-point quadrature are absent; select the
corresponding CPU task instead.
